\documentclass[a4paper]{article}
\usepackage[fontset=fandol]{ctex}
\usepackage{amsmath, amssymb}
\usepackage{graphicx}
\usepackage[margin=2.5cm]{geometry}
\usepackage[most]{tcolorbox}
\usepackage{etoolbox}
\usepackage{listings}
\usepackage{hyperref}
\usepackage{booktabs}
\usepackage{subcaption}
\usepackage{float}
\usepackage{tikz}
\IfFileExists{pgfplots.sty}{\usepackage{pgfplots}\pgfplotsset{compat=1.18}}{}
\newtcolorbox{knowledgebox}[1]{enhanced, colback=blue!5!white, colframe=blue!75!black, colbacktitle=blue!75!black, coltitle=white, fonttitle=\bfseries, title=#1, attach boxed title to top left={yshift=-2mm, xshift=2mm}, boxrule=1pt, sharp corners}
\newtcolorbox{importantbox}[1]{enhanced, colback=yellow!10!white, colframe=yellow!80!black, colbacktitle=yellow!80!black, coltitle=black, fonttitle=\bfseries, title=#1, sharp corners}
\newtcolorbox{warningbox}[1]{enhanced, colback=red!5!white, colframe=red!75!black, colbacktitle=red!75!black, coltitle=white, fonttitle=\bfseries, title=#1, sharp corners}
\lstset{language=Python,basicstyle=\ttfamily\small,keywordstyle=\color{blue},stringstyle=\color{red!60!black},commentstyle=\color{green!60!black},breaklines=true,frame=single,numbers=left,numberstyle=\tiny\color{gray},captionpos=b,extendedchars=false}
\newcommand{\notetitle}{Speech Recognition Formulations}
\newcommand{\noteauthors}{Codex 基于课程页、公开视频字幕、coverage ledger 与补充官方材料整理}
\newcommand{\notedate}{\today}
\newcommand{\videochannel}{WAVLab}
\newcommand{\videopublishdate}{2023-09-06}
\newcommand{\videoduration}{unknown}
\newcommand{\videourl}{https://www.youtube.com/watch?v=Gjc0mSTsoms}
\newcommand{\videocoverpath}{cover.jpg}
\begin{document}
\begin{titlepage}
\centering
{\Large 课程笔记\par}
\vspace{1.2cm}
{\huge\bfseries \notetitle\par}
\vspace{0.8cm}
{\large \noteauthors\par}
\vspace{0.3cm}
{\large \notedate\par}
\vspace{1.2cm}
\ifdefempty{\videocoverpath}{{\small 公开视频封面不可用。\par}}{\includegraphics[width=0.82\textwidth,height=0.45\textheight,keepaspectratio]{\videocoverpath}\par}
\vfill
\begin{tcolorbox}[width=0.9\textwidth, colback=black!2!white, colframe=black!60, sharp corners]
\textbf{视频作者/频道}：\videochannel\par
\textbf{发布日期}：\videopublishdate\par
\textbf{视频时长}：\videoduration\par
\textbf{视频链接}：\href{\videourl}{\nolinkurl{\videourl}}
\end{tcolorbox}
\end{titlepage}
\tableofcontents
\newpage
\section{3.0 课程主题总览}
根据课程页主题、公开视频字幕以及本讲的 coverage ledger，3.0 课程主题总览 这一部分主要围绕 Speech Recognition Formulations | Probabilistic rules | From Bayes decision theory to HMM + n-gram, CTC, RNN-T, and attention 展开。
写作上保持 coverage-first 原则：先把老师在这一阶段真正讲过的问题设定、模型部件、训练或推理逻辑讲清，再压缩成适合复习的结构。

\begin{figure}[H]
\centering
\includegraphics[width=0.92\textwidth]{figures/coverage_map.png}
\caption{本讲 coverage map：按 coverage\_units.jsonl 归纳出的核心主题与章节映射。}
\end{figure}

\begin{figure}[H]
\centering
\includegraphics[width=0.92\textwidth]{figures/timeline.png}
\caption{本讲 timeline：根据字幕时间范围归纳的主题推进顺序。}
\end{figure}

\subsection{Speech Recognition Formulations}
从证据位置看，这一部分对应的课程证据区间是 2023-09-06。课程在这里强调的核心内容可以概括为：Speech Recognition Formulations | Probabilistic rules | From Bayes decision theory to HMM + n-gram, CTC, RNN-T, and attention。
这一步把识别问题写成一个清晰的概率分解，也奠定了后面 HMM、n-gram、以及更现代端到端模型之间的比较框架。

在本章结构中，这一单元被并入“3.0 课程主题总览”，目的是把同一阶段的定义、机制和权衡集中讲清。
\[
\hat{W} = \arg\max_W P(W \mid X) = \arg\max_W P(X \mid W) P(W)
\]

符号说明如下：
\begin{itemize}
\item $X$：观测到的语音特征序列；
\item $W$：候选词序列；
\item $P(X \mid W)$：声学模型或观测模型；
\item $P(W)$：语言模型。
\end{itemize}

\subsection{本章小结}
3.0 课程主题总览 这一部分把 Speech Recognition Formulations 中相邻的教学单元串成了一条连续叙事：既交代了这一阶段要解决的问题，也交代了方法、推理或工程权衡为什么要这样组织。

\section{3.1 概率化建模基础（probabilistic rules）}
根据课程页主题、公开视频字幕以及本讲的 coverage ledger，3.1 概率化建模基础（probabilistic rules） 这一部分主要围绕 Probabilistic rules | Sequence modeling setup | Bayes decision theory 展开。
写作上保持 coverage-first 原则：先把老师在这一阶段真正讲过的问题设定、模型部件、训练或推理逻辑讲清，再压缩成适合复习的结构。

\subsection{Probabilistic rules}
从证据位置看，这一部分对应的课程证据区间是 00:00:00,240 -- 00:06:38,960。课程在这里强调的核心内容可以概括为：Probabilistic rules | Sequence modeling setup | Bayes decision theory。
这一步把识别问题写成一个清晰的概率分解，也奠定了后面 HMM、n-gram、以及更现代端到端模型之间的比较框架。

在本章结构中，这一单元被并入“3.1 概率化建模基础（probabilistic rules）”，目的是把同一阶段的定义、机制和权衡集中讲清。
\[
\hat{W} = \arg\max_W P(W \mid X) = \arg\max_W P(X \mid W) P(W)
\]

符号说明如下：
\begin{itemize}
\item $X$：观测到的语音特征序列；
\item $W$：候选词序列；
\item $P(X \mid W)$：声学模型或观测模型；
\item $P(W)$：语言模型。
\end{itemize}

\subsection{本章小结}
3.1 概率化建模基础（probabilistic rules） 这一部分把 Speech Recognition Formulations 中相邻的教学单元串成了一条连续叙事：既交代了这一阶段要解决的问题，也交代了方法、推理或工程权衡为什么要这样组织。

\section{3.2 Bayes decision theory 与 noisy channel}
根据课程页主题、公开视频字幕以及本讲的 coverage ledger，3.2 Bayes decision theory 与 noisy channel 这一部分主要围绕 Bayes decision theory | Noisy channel factorization | Chain rule 展开。
写作上保持 coverage-first 原则：先把老师在这一阶段真正讲过的问题设定、模型部件、训练或推理逻辑讲清，再压缩成适合复习的结构。

\subsection{Bayes decision theory}
从证据位置看，这一部分对应的课程证据区间是 00:06:36,639 -- 00:13:35,199。课程在这里强调的核心内容可以概括为：Bayes decision theory | Noisy channel factorization | Chain rule。
这一步把识别问题写成一个清晰的概率分解，也奠定了后面 HMM、n-gram、以及更现代端到端模型之间的比较框架。

在本章结构中，这一单元被并入“3.2 Bayes decision theory 与 noisy channel”，目的是把同一阶段的定义、机制和权衡集中讲清。
\[
\hat{W} = \arg\max_W P(W \mid X) = \arg\max_W P(X \mid W) P(W)
\]

符号说明如下：
\begin{itemize}
\item $X$：观测到的语音特征序列；
\item $W$：候选词序列；
\item $P(X \mid W)$：声学模型或观测模型；
\item $P(W)$：语言模型。
\end{itemize}

\subsection{本章小结}
3.2 Bayes decision theory 与 noisy channel 这一部分把 Speech Recognition Formulations 中相邻的教学单元串成了一条连续叙事：既交代了这一阶段要解决的问题，也交代了方法、推理或工程权衡为什么要这样组织。

\section{3.3 HMM + n-gram baseline}
根据课程页主题、公开视频字幕以及本讲的 coverage ledger，3.3 HMM + n-gram baseline 这一部分主要围绕 HMM + n-gram baseline | Tractable decomposition | Sequence probability 展开。
写作上保持 coverage-first 原则：先把老师在这一阶段真正讲过的问题设定、模型部件、训练或推理逻辑讲清，再压缩成适合复习的结构。

\subsection{HMM + n-gram baseline}
从证据位置看，这一部分对应的课程证据区间是 00:13:32,240 -- 00:20:11,640。课程在这里强调的核心内容可以概括为：HMM + n-gram baseline | Tractable decomposition | Sequence probability。
这里的重点不是把 HMM 当作过时历史，而是把它当作对齐、状态建模与动态规划最清晰的教学载体。

在本章结构中，这一单元被并入“3.3 HMM + n-gram baseline”，目的是把同一阶段的定义、机制和权衡集中讲清。
\[
P(w_1,\dots,w_T) = \prod_{t=1}^{T} P(w_t \mid w_1,\dots,w_{t-1}) \approx \prod_{t=1}^{T} P(w_t \mid w_{t-n+1},\dots,w_{t-1})
\]

符号说明如下：
\begin{itemize}
\item $w_t$：第 $t$ 个词；
\item $T$：句子长度；
\item $n$：$n$-gram 的上下文窗口长度。
\end{itemize}

\subsection{本章小结}
3.3 HMM + n-gram baseline 这一部分把 Speech Recognition Formulations 中相邻的教学单元串成了一条连续叙事：既交代了这一阶段要解决的问题，也交代了方法、推理或工程权衡为什么要这样组织。

\section{3.4 CTC（Connectionist Temporal Classification）}
根据课程页主题、公开视频字幕以及本讲的 coverage ledger，3.4 CTC（Connectionist Temporal Classification） 这一部分主要围绕 CTC | Monotonic alignments | Path summation 展开。
写作上保持 coverage-first 原则：先把老师在这一阶段真正讲过的问题设定、模型部件、训练或推理逻辑讲清，再压缩成适合复习的结构。

\subsection{CTC}
从证据位置看，这一部分对应的课程证据区间是 00:20:09,799 -- 00:27:16,720。课程在这里强调的核心内容可以概括为：CTC | Monotonic alignments | Path summation。
CTC 的关键在于保持单调对齐假设，同时把对齐路径视作潜变量并在训练时进行求和。

在本章结构中，这一单元被并入“3.4 CTC（Connectionist Temporal Classification）”，目的是把同一阶段的定义、机制和权衡集中讲清。
\[
P(Y \mid X) = \sum_{\pi \in \mathcal{B}^{-1}(Y)} \prod_{t=1}^{T} P(\pi_t \mid X)
\]

符号说明如下：
\begin{itemize}
\item $X$：输入语音特征序列；
\item $Y$：输出标签序列；
\item $\pi$：带有 blank 的对齐路径；
\item $\mathcal{B}$：把路径压缩为最终标签序列的映射。
\end{itemize}

\subsection{本章小结}
3.4 CTC（Connectionist Temporal Classification） 这一部分把 Speech Recognition Formulations 中相邻的教学单元串成了一条连续叙事：既交代了这一阶段要解决的问题，也交代了方法、推理或工程权衡为什么要这样组织。

\section{3.5 RNN-T / transducer}
根据课程页主题、公开视频字幕以及本讲的 coverage ledger，3.5 RNN-T / transducer 这一部分主要围绕 RNN-T | Prediction network | Transducer objective 展开。
写作上保持 coverage-first 原则：先把老师在这一阶段真正讲过的问题设定、模型部件、训练或推理逻辑讲清，再压缩成适合复习的结构。

\subsection{RNN-T}
从证据位置看，这一部分对应的课程证据区间是 00:27:14,919 -- 00:33:34,960。课程在这里强调的核心内容可以概括为：RNN-T | Prediction network | Transducer objective。
课程把 RNN-T 放在 streaming ASR 的语境里讨论，强调它保留了单调约束，却让模型可以更灵活地联合时间和输出历史。

在本章结构中，这一单元被并入“3.5 RNN-T / transducer”，目的是把同一阶段的定义、机制和权衡集中讲清。
\[
P(Y \mid X) = \sum_{\pi \in \mathcal{A}(X,Y)} \prod_{(t,u)\in\pi} P(k_{t,u} \mid h_t, g_u)
\]

符号说明如下：
\begin{itemize}
\item $X$：输入语音特征序列；
\item $Y$：目标标签序列；
\item $\mathcal{A}(X,Y)$：满足转导约束的对齐集合；
\item $h_t$：encoder 在时间步 $t$ 的表示；
\item $g_u$：prediction network 在输出步 $u$ 的表示。
\end{itemize}

\subsection{本章小结}
3.5 RNN-T / transducer 这一部分把 Speech Recognition Formulations 中相邻的教学单元串成了一条连续叙事：既交代了这一阶段要解决的问题，也交代了方法、推理或工程权衡为什么要这样组织。

\section{3.6 attention-based seq2seq}
根据课程页主题、公开视频字幕以及本讲的 coverage ledger，3.6 attention-based seq2seq 这一部分主要围绕 Attention-based seq2seq | Soft alignments | Decoding tradeoffs 展开。
写作上保持 coverage-first 原则：先把老师在这一阶段真正讲过的问题设定、模型部件、训练或推理逻辑讲清，再压缩成适合复习的结构。

\subsection{Attention-based seq2seq}
从证据位置看，这一部分对应的课程证据区间是 00:33:32,000 -- 00:40:25,040。课程在这里强调的核心内容可以概括为：Attention-based seq2seq | Soft alignments | Decoding tradeoffs。
与单调模型不同，attention encoder-decoder 更强调全局条件化能力，但也因此引入了对齐自由度和部署层面的权衡。

在本章结构中，这一单元被并入“3.6 attention-based seq2seq”，目的是把同一阶段的定义、机制和权衡集中讲清。

\subsection{本章小结}
3.6 attention-based seq2seq 这一部分把 Speech Recognition Formulations 中相邻的教学单元串成了一条连续叙事：既交代了这一阶段要解决的问题，也交代了方法、推理或工程权衡为什么要这样组织。

\section{3.7 统一视角与 tradeoffs}
根据课程页主题、公开视频字幕以及本讲的 coverage ledger，3.7 统一视角与 tradeoffs 这一部分主要围绕 Unified view | Model tradeoffs | Low-resource and practical remarks 展开。
写作上保持 coverage-first 原则：先把老师在这一阶段真正讲过的问题设定、模型部件、训练或推理逻辑讲清，再压缩成适合复习的结构。

\subsection{Unified view}
从证据位置看，这一部分对应的课程证据区间是 00:40:22,880 -- 00:46:44,000。课程在这里强调的核心内容可以概括为：Unified view | Model tradeoffs | Low-resource and practical remarks。
从课程组织上看，这一部分承担的是把局部概念嵌回整条识别流水线中的作用，避免读者把公式或模块孤立地记忆。

在本章结构中，这一单元被并入“3.7 统一视角与 tradeoffs”，目的是把同一阶段的定义、机制和权衡集中讲清。

\subsection{本章小结}
3.7 统一视角与 tradeoffs 这一部分把 Speech Recognition Formulations 中相邻的教学单元串成了一条连续叙事：既交代了这一阶段要解决的问题，也交代了方法、推理或工程权衡为什么要这样组织。

\section{总结与延伸}
Speech Recognition Formulations 这一讲在整门《Speech Recognition and Understanding》课程中的位置，是把 Speech Recognition Formulations 讲成一个可复用的建模模块。
从教材化角度看，本讲最重要的不是记住单一术语，而是理解它如何嵌入完整的 automatic speech recognition pipeline：输入语音如何被表示、模型如何被训练、解码如何得到最终转写，以及这些设计分别带来什么假设和权衡。
由于公开源以课程页和公开视频字幕为主、缺少官方 slide PDF，本章在正文里尽量把术语、方法和工程语境解释完整，并把缺失项留在 omission 和 manifest 里，而不是凭空补写。
\end{document}
